\section*{第 \ref{ch:g11:work-of-force} 章——力做的功}

\begin{solution}{exo:g11:work-of-force:1}
$W = 40 \times 50 \times \cos\theta$：$+\qty{2000}{J}$；
$+\qty{1000}{J}$；$0$；$-\qty{1000}{J}$。
\end{solution}

\begin{solution}{exo:g11:work-of-force:2}
$mgh = 12 \times 9.81 \times 2.5 = \qty{294}{J}$：下降时，
$+\qty{294}{J}$（\omterm{def:g10:universal-gravitation:weight}{重力}沿下降方向）；上升时，$-\qty{294}{J}$。
\end{solution}

\begin{solution}{exo:g11:work-of-force:3}
(a) 电动机做功；(b) 阻力做功；(c) 零（垂直）；(d) 阻力做功；
(e) 零——绳子的\omterm{def:g10:inertia:inventory}{张力}始终与摆球
运动方向垂直。
\end{solution}

\begin{solution}{exo:g11:work-of-force:4}
匀速时，提升力 $= mg = \qty{5.89e3}{N}$，因此
$W_{\text{提升}} = 5886 \times 25 \approx +\qty{147}{kJ}$，且
$W(\vect P) = -\qty{147}{kJ}$。总功为零：匀速运动时
各力平衡（见\cref{ch:g11:forces-and-motion}），因此它们所做的功
\omterm{def:g10:inertia:compensate}{相互抵消}。
\end{solution}

\begin{solution}{exo:g11:work-of-force:5}
$P = \num{9.0e4}/60 = \qty{1.5}{kW}$；$1500/736 \approx \qty{2.0}{hp}$。
\end{solution}

\begin{solution}{exo:g11:work-of-force:6}
绳子拉力：$120 \times 200 \times \cos\qty{30}{\degree} \approx
+\qty{20.8}{kJ}$；\omterm{def:g10:inertia:inventory}{摩擦力}：$-45 \times 200 = -\qty{9.0}{kJ}$；\omterm{def:g10:universal-gravitation:weight}{重力}
和\omterm{def:g10:inertia:inventory}{法向力}与位移垂直，做功均为 $0$。总功 $+\qty{11.8}{kJ}$：
雪橇的速率增大。
\end{solution}

\begin{solution}{exo:g11:work-of-force:7}
$v = 25/3.6 = \qty{6.94}{m/s}$；$P = Fv = 18 \times 6.94 \approx
\qty{125}{W}$。一小时内，$W = P\,\Delta t = 125 \times 3600 =
\qty{450}{kJ}$（亦可写为 $18 \times \qty{25000}{m}$）。
\end{solution}

\begin{solution}{exo:g11:work-of-force:8}
两条路线上的\omterm{def:g10:universal-gravitation:weight}{重力}做功：$+65 \times 9.81 \times 120 \approx
+\qty{76.5}{kJ}$（仅下降段做功）。\omterm{def:g10:inertia:inventory}{摩擦力}：$-30 \times 800 =
-\qty{24}{kJ}$（斜坡上），$-30 \times 1500 = -\qty{45}{kJ}$（
\omterm{def:g10:universal-gravitation:satellite}{轨道}上）。只有\omterm{def:g10:inertia:inventory}{摩擦力}做功与路径有关。
\end{solution}

\begin{solution}{exo:g11:work-of-force:9}
上升高度 $h = 4.0\sin\qty{25}{\degree} = \qty{1.69}{m}$：
$W(\vect P) = -20 \times 9.81 \times 1.69 \approx -\qty{332}{J}$；
$W_f = -30 \times 4.0 = -\qty{120}{J}$。匀速运动时总功
为零，因此 $W_{\text{推力}} = +\qty{452}{J}$，且
$F = 452/4.0 \approx \qty{113}{N}$。
\end{solution}

\begin{solution}{exo:g11:work-of-force:10}
$v = 130/3.6 = \qty{36.1}{m/s}$；匀速运动，因此
$F = P/v = \num{5.0e4}/36.1 \approx \qty{1.4}{kN}$。
\end{solution}

\begin{solution}{exo:g11:work-of-force:11}
$P = Fv\cos\theta = 800 \times 1.0 \times \cos\qty{20}{\degree}
\approx \qty{752}{W}$；以\omterm{def:g11:work-of-force:power}{马力}计，$752/736 \approx \num{1.02}$——
几乎恰好是一匹强壮的马能持续输出的\omterm{def:g11:work-of-force:power}{功率}，
\omterm{def:g11:work-of-force:power}{瓦特}正是以此比较来确定这一\omterm{def:g10:orders-of-magnitude:unit}{单位}的。
\end{solution}

\begin{solution}{exo:g11:work-of-force:12}
长度 $L = h/\text{坡度}$：\qty{1500}{m} 和 \qty{5000}{m}。力
$F = mg \times \text{坡度}$：$1200 \times 9.81 \times 0.20 \approx
\qty{2.35}{kN}$ 和 $\qty{0.71}{kN}$。功 $FL$：$2354 \times 1500
\approx \qty{3.53}{MJ}$ 和 $706 \times 5000 \approx \qty{3.53}{MJ}$
——二者均等于 $mgh$：道路只是以长度换取力，功永远不变。
\end{solution}

\begin{solution}{exo:g11:work-of-force:13}
$F = mg = \qty{2.45e3}{N}$；$W = 2453 \times 15 \approx \qty{36.8}{kJ}$；
$\Delta t = 15/0.40 = \qty{37.5}{s}$；
$P = 36800/37.5 \approx \qty{981}{W} = 2453 \times 0.40$（与 $Fv$ 结果一致）；
输入电\omterm{def:g11:work-of-force:power}{功率}：$981/0.70 \approx \qty{1.4}{kW}$。
\end{solution}

\begin{solution}{exo:g11:work-of-force:14}
对角线：$-140 \times 6.0 = -\qty{840}{J}$；墙边：
$-140 \times 8.4 \approx -\qty{1.18}{kJ}$；往返一趟：
$-140 \times 12 = -\qty{1.68}{kJ}$。往返一趟时，\omterm{def:g10:universal-gravitation:weight}{重力}做功为零
（没有净下降）。\omterm{def:g10:inertia:inventory}{摩擦力}做的功取决于路径，且永远
不会收回；\omterm{def:g10:universal-gravitation:weight}{重力}做的功只取决于端点
（见\cref{rem:g11:work-of-force:divide}）。
\end{solution}

\begin{solution}{exo:g11:work-of-force:15}
\omterm{def:g10:universal-gravitation:weight}{重力}：$80 \times 9.81 \times 600 \approx \qty{471}{kJ}$；\omterm{def:g10:inertia:inventory}{摩擦力}：
$14 \times 12000 = \qty{168}{kJ}$；总功 $\qty{639}{kJ}$。
$v = \qty{4.17}{m/s}$，因此 $\Delta t = 12000/4.17 = \qty{2880}{s} =
\qty{48}{min}$，$P = \num{6.39e5}/2880 \approx \qty{222}{W} \approx
\qty{0.30}{hp}$——一位体力很好的业余选手，相当于一匹马的三分之一。
\end{solution}

\begin{solution}{pb:g11:work-of-force:1}
\textbf{1.} $W = 110 \times 8.0 = +\qty{880}{J}$。

\textbf{2.} \omterm{def:g10:inertia:inventory}{摩擦力}：$-95 \times 8.0 = -\qty{760}{J}$。\omterm{def:g10:universal-gravitation:weight}{重力}和
\omterm{def:g10:inertia:inventory}{法向力}与运动方向垂直，做功均为零。

\textbf{3.} 总功 $+\qty{120}{J}$：\omterm{def:g11:forces-and-motion:netforce}{合力}做的正功，箱子的速率增大。
用 \qty{95}{N} 的推力与\omterm{def:g10:inertia:inventory}{摩擦力}平衡——\omterm{def:g11:forces-and-motion:netforce}{合力}为零，速率
不变，总功为零。

\textbf{4.} 对角线：$-95 \times 10.0 = -\qty{950}{J}$；沿墙：
$-95 \times 14.0 = -\qty{1.33}{kJ}$：\omterm{def:g10:inertia:inventory}{摩擦力}按每米“收费”。

\textbf{5.} \omterm{def:g10:inertia:inventory}{摩擦力}：$-95 \times 20.0 = -\qty{1.9}{kJ}$；\omterm{def:g10:universal-gravitation:weight}{重力}：
$0$（净下降为零）。只有\omterm{def:g10:universal-gravitation:weight}{重力}会“退还”。

\textbf{6.} $W(\vect P) = -30 \times 9.81 \times 9.0 =
-\qty{2.65}{kJ}$；搬运者至少须提供 $+\qty{2.65}{kJ}$ 的功。

\textbf{7.} $(70 + 30) \times 9.81 \times 9.0 = \qty{8.83}{kJ}$；
箱子所占份额为 $2.65/8.83 \approx 30\%$——搬运
主要是在运送人自身。

\textbf{8.} $\sin\alpha = 9.0/30 = 0.30$:
$F = mg\sin\alpha + f = 88.3 + 40 = \qty{128}{N}$;
$W = 128 \times 30 = \qty{3.85}{kJ}$ ($= 2.65 + 1.2$).

\textbf{9.} 竖直吊起：$F = mg = \qty{294}{N}$，
$W = \qty{2.65}{kJ}$。斜面以更多的功为代价换得较小的力（$128$ 对
比 \qty{294}{N}），即\omterm{def:g10:inertia:inventory}{摩擦力}收取的 \qty{1.2}{kJ}
“过路费”。

\textbf{10.} $\Delta t = 9.0/0.50 = \qty{18}{s}$;
$P = 2649/18 \approx \qty{147}{W}$.

\textbf{11.} $P = Fv = 294 \times 0.50 = \qty{147}{W}$——结果相同。
\omterm{def:g10:energy-conservation:forms}{电能}输入：$147/0.60 \approx \qty{245}{W}$。

\textbf{12.} $1000 \times 0.24 = \qty{240}{m}$ 与
$8000 \times 0.030 = \qty{240}{m}$：到达同一高度。

\textbf{13.} $F = mg \times \text{坡度}$：车道
$3500 \times 9.81 \times 0.24 \approx \qty{8.2}{kN}$；公路
$\approx \qty{1.0}{kN}$.

\textbf{14.} $8240 \times 1000 \approx \qty{8.24}{MJ}$ 与
$1030 \times 8000 \approx \qty{8.24}{MJ}$——两者均为 $= mgh = 3500
\times 9.81 \times 240$.

\textbf{15.} 滚动阻力：$400 \times 1000 = \qty{0.40}{MJ}$
相比 $400 \times 8000 = \qty{3.2}{MJ}$：平缓道路的舒适
需额外多做 \qty{2.8}{MJ} 的摩擦功。

\textbf{16.} $v = \qty{10}{m/s}$。道路：
$(1030 + 400) \times 10 = \qty{14.3}{kW} \approx \qty{19}{hp}$；车道：
$(8240 + 400) \times 10 = \qty{86.4}{kW} \approx \qty{117}{hp}$ ---
超出满载货车的能力。道路蜿蜒盘绕，使\omterm{def:g11:work-of-force:power}{功率}适中的发动机（及轮胎）
也能爬坡：做功相同，只是摊得更薄。

\textbf{17.} 每箱：$880 + 2649 = \qty{3.53}{kJ}$；$60$ 箱共计
$60 \times 3529 \approx \qty{212}{kJ}$.

\textbf{18.} 楼梯：$20 \times 70 \times 9.81 \times 9.0 \approx
\qty{124}{kJ}$。全天合计 $\approx \qty{336}{kJ} = \qty{0.34}{MJ}$。

\textbf{19.} $P = \num{3.36e5}/28800 \approx \qty{12}{W}$。水壶：
$\num{3.36e5}/2000 = \qty{168}{s}$——烧水不到三分钟，相当于
全天的劳动。

\textbf{20.} \omterm{def:g11:work-of-force:work}{焦耳}只计量力\emph{沿}运动方向所做的功——
推墙或托着箱子都不计功，无论多么费力
（即使机械功为零，肌肉仍会消耗\omterm{def:g10:energy-conservation:forms}{化学能}）。由于
人的\omterm{def:g11:work-of-force:power}{功率}输出仅有几十瓦，搬运者便让机器来改变
力的方式——绞车、斜面、之字形路线——而 $mgh$ 这笔账始终
不变。
\end{solution}
